\documentclass[letterpaper,11pt]{article}

\usepackage{latexsym}
\usepackage[empty]{fullpage}
\usepackage{titlesec}
\usepackage{marvosym}
\usepackage[usenames,dvipsnames]{color}
\usepackage{verbatim}
\usepackage{enumitem}
\usepackage[hidelinks]{hyperref}
\usepackage{fancyhdr}
\usepackage{tabularx}
\usepackage{xcolor}
\usepackage{fontawesome5}
\usepackage[most]{tcolorbox}
\usepackage{amssymb}

% -------------------- RUSSIAN / CYRILLIC SUPPORT --------------------
% These three lines were added so Cyrillic text renders with pdfLaTeX.
\usepackage[utf8]{inputenc}
\usepackage[T1,T2A]{fontenc}
\usepackage[english,russian]{babel}

\input{glyphtounicode}

% -------------------- FONT OPTIONS --------------------
% sans-serif
% \usepackage[sfdefault]{roboto}
% \usepackage[sfdefault]{noto-sans}
\usepackage{CormorantGaramond}
% NOTE: Cormorant Garamond may not ship Cyrillic glyphs. If your Russian text
% renders in a mismatched fallback font (or you see "Font shape ... undefined"
% warnings), comment out the line above and uncomment the line below instead:
% \usepackage{paratype}   % PT Serif/Sans — designed for Cyrillic, works with pdfLaTeX
% serif
% \usepackage{charter}

\pagestyle{fancy}
\fancyhf{} % clear all header and footer fields
\fancyfoot{}
\renewcommand{\headrulewidth}{0pt}
\renewcommand{\footrulewidth}{0pt}

% Adjust margins
\addtolength{\oddsidemargin}{-0.5in}
\addtolength{\evensidemargin}{-0.5in}
\addtolength{\textwidth}{1in}
\addtolength{\topmargin}{-1in} % Default was -.5in
\addtolength{\textheight}{1.0in}

\urlstyle{same}

\raggedbottom
\raggedright
\setlength{\tabcolsep}{0in}

% Section formatting
\titleformat{\section}{
	\vspace{-5pt}\scshape\raggedright\large
}{}{0em}{}[\color{black}\titlerule \vspace{-5pt}]

% Subsection formatting
\titleformat{\subsection}{
	\vspace{-4pt}\scshape\raggedright\large
}{\hspace{-.15in}}{0em}{}[\color{black}\vspace{-8pt}]

% Ensure that generate pdf is machine readable/ATS parsable
\pdfgentounicode=1

% -------------------- CUSTOM COMMANDS --------------------
\newcommand{\resumeItem}[1]{
	\item\small{
		{#1 \vspace{-2pt}}
	}
}

\newcommand{\resumeSubheading}[4]{
	\vspace{-2pt}\item
	\begin{tabular*}{0.97\textwidth}[t]{l@{\extracolsep{\fill}}r}
		\textbf{#1} & #2 \\
		\textit{\small#3} & \textit{\small #4} \\
	\end{tabular*}\vspace{-7pt}
}

\newcommand{\resumeSubSubheading}[2]{
	\item
	\begin{tabular*}{0.97\textwidth}{l@{\extracolsep{\fill}}r}
		\textit{\small#1} & \textit{\small #2} \\
	\end{tabular*}\vspace{-7pt}
}

\newcommand{\resumeProjectHeading}[2]{
	\item
	\begin{tabular*}{0.97\textwidth}{l@{\extracolsep{\fill}}r}
		\small#1 & #2 \\
	\end{tabular*}\vspace{-7pt}
}

\newcommand{\resumeSubItem}[1]{\resumeItem{#1}\vspace{-4pt}}
\newcommand{\resumeSubHeadingListStart}{\begin{itemize}[leftmargin=0.15in, label={}]}
		\newcommand{\resumeSubHeadingListEnd}{\end{itemize}}
\newcommand{\resumeItemListStart}{\begin{itemize}}
		\newcommand{\resumeItemListEnd}{\end{itemize}\vspace{-5pt}}

\renewcommand\labelitemii{$\vcenter{\hbox{\tiny$\bullet$}}$}

\setlength{\footskip}{4.08003pt}

\newtcbox{\tooltag}[1][black]{on line,
	enhanced,          % needed for arc
	arc=4pt,           % small rounding
	colback=white,     % light fill
	colframe=#1!60,     % stronger border
	boxrule=0.2pt,     % thin border
	boxsep=0.1pt,      % padding
	left=5pt, right=5pt, top=-2pt, bottom=-1pt,
	fontupper=\strut\color{black}\rmfamily\fontsize{8}{1}\selectfont
}

% -------------------- START OF DOCUMENT --------------------
\begin{document}

% -------------------- HEADING--------------------
\vspace{10pt}
\begin{center}
\end{center}

\begin{center}
	\textbf{\Large \scshape Владислав Балабаев} \\ \vspace{8pt}
	\small
	\faIcon{github}
	\href{https://github.com/vbalab}{\underline{github.com/vbalab}} $  $
	\faIcon{telegram}
	\href{https://t.me/vbalab}{\underline{t.me/vbalab}} $  $
	\faIcon{envelope}
	\href{mailto:vv.balabaev@gmail.com}
	{\underline{vv.balabaev@gmail.com}} $  $
	% \faIcon{map-marked-alt} Москва, Россия
\end{center}

% -------------------- EDUCATION --------------------
\section{Образование}
\resumeSubHeadingListStart

\resumeSubheading
{Школа анализа данных Яндекса (ШАД)}{2024 -- 2026}
{Инфраструктура больших данных}{Средний балл: 4.5/5.0}

\resumeSubheading
{Китайский университет Гонконга, Шэньчжэнь}{2024 -- 2025}
{Программа обмена}{Средний балл: 4.8/5.0}

\resumeSubheading
{Российская экономическая школа (РЭШ)}{2023 -- 2026}
{Магистратура по экономике}{Средний балл: 4.8/5.3}

\resumeSubheading
{Новосибирский государственный университет}{2019 -- 2023}
{Бакалавр экономики, прикладная эконометрика}{Средний балл: 4.7/5.0}

\resumeSubHeadingListEnd

% -------------------- WORK EXPERIENCE --------------------
\section{Опыт работы}
\resumeSubHeadingListStart

\resumeProjectHeading
{\textbf{ML-инженер, \href{https://autonomy.yandex.ru}{\underline{Яндекс, Автономный Транспорт, ML R\&D}}} \vspace{8pt}}{2026-01 -- 2026-05}
{\small{Исследовал и прототипировал архитектурные улучшения для онлайн-построения HD-карт.\\Ключевой вклад: модули, учитывающие топологию, альтернативные геометрические функции потерь, учет времени}}

\resumeProjectHeading
{\textbf{Дата-сайентист, \href{https://www.vtb.ru}{\underline{Банк ВТБ}}} \vspace{8pt}}{2025-06 -- 2026-01}
{\small{Предиктивное моделирование, разработка пайплайна CLTV}}

\resumeProjectHeading
{\textbf{Дата-сайентист, \href{https://okkam.group/}{\underline{Okkam}}} \vspace{8pt}}{2024-06 -- 2024-08}
{\small{Вел проект по marketing mix modeling, проводил продвинутый ad-hoc анализ данных}}

\resumeProjectHeading
{\textbf{Дата-сайентист, \href{https://catalog.cft.ru/}{\underline{Центр финансовых технологий (ЦФТ)}}} \vspace{8pt}}{2023-01 -- 2023-09}
{\small{Поддерживал сквозной пайплайн кредитного скоринга, оптимизировал инженерию признаков и выбор моделей. \\ Выявлял паттерны мошенничества с помощью статистического анализа и проверки гипотез}}

\resumeProjectHeading
{\textbf{Дата-инженер, стажировка,  \href{https://www.sberbank.com/index}{\underline{СБЕР}}}
	\vspace{8pt}}{2022-09 -- 2022-12}
{\small{Спроектировал и выпустил библиотеку Greenplum-PySpark для DS-команды. \\
Автоматизировал мониторинг витрин данных с помощью SQL- и PySpark-пайплайнов}}

\resumeProjectHeading
{\textbf{Научный сотрудник, \href{https://www.nsu.ru/n/media/news/obrazovanie/v-ngu-sozdana-novaya-laboratoriya-analiziruyushchaya-otraslevye-rynki/}{\underline{Лаборатория эмпирического анализа отраслевых рынков}}}
	\vspace{8pt}}{2021-09 -- 2024-08}
{\small{Собрал и проанализировал более 3,5 млн объявлений о продаже автомобилей в России для моделирования оптимальной налоговой политики. \\ Моделировал спрос для розничной сети «Семья» с помощью эконометрики и ML}}

% -------------------- PROJECTS --------------------
\section{Проекты}

\resumeProjectHeading
{\textbf{\href{https://t.me/NesRandomCoffeeBot}{\underline{NESpresso}} -- официальный Telegram-бот выпускников РЭШ} \vspace{8pt}}{}
{\tooltag{OpenSearch} \tooltag{RAG} \tooltag{Aiogram} \tooltag{FastAPI} \tooltag{Postgres} \tooltag{Docker} \\
	\small{Объединяет более 3500 выпускников РЭШ через встречи и умный поиск по запросу}}

% -------------------- PROFESSIONAL SKILLS --------------------
\section{Профессиональные навыки}
\begin{itemize}[leftmargin=0.15in, label={}]
	\small{\item	      {
	      \vspace{3pt}
	      $\circ$ \textbf{ML}{: классический ML, причинно-следственное и эконометрическое моделирование, NLP, CV, генеративное моделирование}
	      \vspace{2pt}

	      $\circ$ \textbf{C{\Large ++}}{: ООП и паттерны проектирования, многопоточность, сетевое программирование, метапрограммирование}
	      \vspace{2pt}

	      $\circ$ \textbf{Python}{: fastapi, sqlalchemy, asyncio, pyspark, torch, catboost, plotly}
	      \vspace{2pt}

	      $\circ$ \textbf{Базы данных}{: SQLAlchemy, Postgres, Greenplum, экосистема Hadoop}
	      \vspace{2pt}

	      $\circ$ \textbf{Инструменты}{: Linux/Unix, Docker, Git, Apache Spark, Redis, MLflow}
	      \vspace{3pt}
	      }}
\end{itemize}

\resumeSubHeadingListEnd

\end{document}